\begin{lemma}[Scheide three departures and replacement data (7)]
For every defective color $\delta$, the boundary identities of color
$\delta$ from $Z_{\gamma(\delta)}$ to the exterior of $V(T)$ are exactly
three.  They can be named $f^0_\delta,f^1_\delta,f^2_\delta$ with distinct
outside ends, and these choices, together with the exact replacement
tree-sequence obtained by deleting the $\gamma(\delta)$ block and adjoining
$f^1_\delta,f^2_\delta$, form a family of
\noderef{ScheideReplacementData}.
\end{lemma}
\begin{proof}
Fix $\delta$ and abbreviate $Z=Z_{\gamma(\delta)}$.  Every vertex of $Z$
is incident with a unique $\delta$-edge: defectiveness plus closedness makes
the color present, and properness makes the incidence unique.  The set $Z$
has odd order, since the globally maximum tree has odd order and exactly two
vertices were deleted.  By \noderef{ScheideDefectiveNoCross}, no
$\delta$-edge joins $Z$ to the deleted pair.  Counting the unique
$\delta$ incidence at each vertex of $Z$ therefore shows that the number
of $\delta$-edges from $Z$ to the exterior of $V(T)$ is odd.  It contains
$f^0_\delta$.

If there were more than three, choose three identities other than
$f^0_\delta$.  The exact switch statement
\noderef{ScheideDefectiveSwitchEffect} says that all three retain color
$\delta$ under $\varphi_\delta$, while that coloring is proper and leaves
exactly $e_0$ uncolored.  Delete the two vertices of the
$\gamma(\delta)$ block and append the three chosen boundary identities and
their outside ends.  Their ends are distinct by properness.  Since
$\delta$ is missing at $v^0_\delta$ after the switch, the appendings satisfy
the Tashkinov rule and yield one more vertex than the global maximum, a
contradiction.

If the odd number were less than three, it would be one, with sole identity
$f^0_\delta$.  By \noderef{ScheideStructuralInternalPair}, an internal
$\alpha_0$-edge joins the deleted pair.  The no-cross assertion excludes
$\delta$ between that pair and $Z$, while closedness excludes
$\alpha_0$ from $V(T)$ to its exterior.  Thus the initial
$(\alpha_0,\delta)$ path avoids the deleted pair.  Consider the distinct
$(\alpha_0,\delta)$ component through one deleted vertex.  Its internal
$\alpha_0$-edge and the no-cross assertion keep it away from $Z$; it uses
no tree edge identity, because the only tree edges incident with the pair
have color $\gamma(\delta)$.  Switching this component preserves
properness by \noderef{TwoColorComponentSwapProper}, fixes the unique
uncolored edge, and preserves every Tashkinov witness.  At least two of the
three or more defective $\delta$ boundary identities must be incident with
the deleted pair.  They become $\alpha_0$-edges, one of which can be
appended in the color missing at $y_0$.  This again exceeds the maximum.
Hence the number is exactly three.

Name the other two identities $f^1_\delta,f^2_\delta$, their ends in $Z$
$v^1_\delta,v^2_\delta$, and their outside ends
$u^1_\delta,u^2_\delta$.  Looplessness gives the displayed endpoint
equations, and properness gives pairwise distinct edge identities and
distinct outside ends.  Delete the named block vertices and block edge
identities from the ordered sequence and append
$f^1_\delta,u^1_\delta,f^2_\delta,u^2_\delta$.  The direct-to-base
attachment fields of \noderef{ScheideReducedConfiguration} show that the
remaining order is a tree-sequence, while the two new edges attach their
outside vertices.  Its vertex and edge images are exactly the finite sets
stated in \noderef{ScheideReplacementData}.  Finite choice over all
defective colors produces the asserted replacement-data family.
\end{proof}
